
\subsection{Описание расчетного блока}

На данном этапе выбираются соответствующие решаемой задаче методы расчета показателей и определения критериев для их интерпретации.

Опираясь на здравый смысл и выполненный в 
%разделе №№ 
рамках 1 этапа НИР анализ источников можно разделить все методы анализа прибрежных территорий на две группы: методы, описывающие влияние со стороны моря, и методы, описывающие состояние суши.
При этом очевидно, что \enquote{морские} методы обладают более высокой специфичностью и приоритетом, так как именно их наличие и определяет \enquote{прибрежность} рассматриваемой территории.

Влияние моря на берег должно учитываться во всех вариантах решаемых в методе задач, а методы анализа рельефа, геологии и исторической динамики береговой линии -- лишь в частных случаях.
Исходя из этого, приоритет внимания в данной работе выделен оценке механического влияния моря на берег и вынесенный в отдельный раздел~\ref{sec:wave_theory} на стр.~\pageref{sec:wave_theory}.

Несмотря на то, что хронологически расчетный модуль идет после блока получения и подготовки геопространственной информации, именно он определяет требования к требуемой структуре и полноте необходимых исходных данных.
Исходя их этого, руководствуясь известным принципом инверсии зависимостей, представляется целесообразным нарушить в описании последовательный порядок изложения методики по выполнению и рассмотреть в первую очередь основной расчетный блок.

% Частные случаи использования данных о состоянии суши рассмотрены в контексте описания конкретных примеров и методов решения подзадач анализа прибрежных территорий и описаны в подразделе №№. 
